\openright\chapter{Samenvatting}

Het landschap van de theoretische vastestoffysica is de afgelopen eeuw drastisch veranderd, waarbij computers nu een centrale rol spelen in wetenschappelijke ontdekkingen. In dit proefschrift passen we de technieken uit informatica toe om kwantum veeldeeltjessimulaties te vereenvoudigen, te versnellen en te verbeteren.

De uitdaging in kwantum veeldeeltjessimulaties ontstaat door een fundamenteel verschil tussen klassieke en kwantumsystemen. In de klassieke mechanica vereist het beschrijven van de beweging van een biljartbal slechts 2 coördinaten (bijvoorbeeld $x$- en $y$-posities, ervan uitgaande dat de bal alleen horizontaal beweegt). Voor 3 ballen hebben we $2+2+2=6$ coördinaten nodig; voor 10 ballen hebben we 20 coördinaten nodig. Het aantal coördinaten groeit lineair met het aantal deeltjes, waardoor het eenvoudig is om duizenden deeltjes op een computer te simuleren.

Kwantummechanica is anders. Als we 2 getallen nodig hebben om de toestand van één kwantumdeeltje te beschrijven, dan vereist het beschrijven van 3 dergelijke deeltjes in het algemeen $2\cdot 2\cdot 2=8$ getallen. Het beschrijven van 10 deeltjes vereist dan $2^{10}=1024$ getallen. Deze exponentiële schaling is fundamenteel voor de kwantum veeldeeltjestheorie en maakt het simuleren van zelfs 100 deeltjes extreem uitdagend, omdat men $1267650600228229401496703205376$ getallen zou moeten opslaan.

Ondanks deze exponentiële schaling zijn er een aantal gevallen waarin we kunnen hopen het probleem te kunnen oplossen. Hoofdstukken~\ref{ch:prb18} en~\ref{ch:mat22} focussen op zwak wisselwerkende deeltjes, waarbij de vereiste computergeheugen terugvalt op dezelfde lineaire schaling als in de klassieke mechanica. Voor een klein aantal sterk wisselwerkende deeltjes kunnen we het probleem exact oplossen. Deze aanpak wordt besproken in hoofdstukken~\ref{ch:ls21},~\ref{ch:paw-atm}, en~\ref{ch:corr-hops}. Tot slot, bieden moderne machine learning methodes benaderende oplossingen voor sterk wisselwerkende gevallen, hoewel ze een analogie van het zogenaamde tekenprobleem tegenkomen. Hoofdstukken~\ref{ch:nc20} en~\ref{ch:cp23} onderzoeken dit probleem en introduceren nieuwe algoritmes om het aan te pakken.

De technieken die in dit werk zijn ontwikkeld, verbeteren onze computationele mogelijkheden in de kwantum veeldeeltjestheorie en tonen aan hoe interdisciplinaire benaderingen uitdagende numerieke problemen kunnen aanpakken.
